\section{Problem Formulation}
\label{sec:problem}

\subsection{Threat Model}

We consider a chosen-plaintext attacker (CPA) with oracle access to an $r$-round block cipher $E^r_K$ keyed by a uniformly random secret key $K$.
The attacker knows a fixed input difference $\Delta P$ and receives ciphertext pairs $(C, C')$ produced from plaintext pairs $(P, P \oplus \Delta P)$.
The attacker does not know the key and has no access to intermediate round states.

\paragraph{Success metric.}
Let $f: \{0,1\}^{2b} \to [0,1]$ be a binary classifier that takes a ciphertext pair and outputs the probability that it was generated from the real distribution (label 1) versus the random distribution (label 0).
We define the \emph{distinguishing advantage} as
\begin{equation}
  \varepsilon = \mathrm{acc}(f) - 0.5,
\end{equation}
where $\mathrm{acc}(f)$ is the classification accuracy on a balanced test set.
A distinguisher is \emph{successful} if $\varepsilon > 0$ with statistical significance ($p < 0.05$ under a one-sample $t$-test against 0.5).

\paragraph{Classification in the $n$-$m$-$T$-$E$ framework.}
Following the taxonomy of Bellini et al.~\citep{bellini2023dbitnet} and Gérault et al.~\citep{gerault2024sok}, our setup is \textbf{2-1-CT-R}: $n=2$ ciphertext inputs, $m=1$ output (binary label), ciphertext-only input type ($T=\text{CT}$), and Gohr's ``real vs.\ random'' labeling ($E=\text{R}$).

\subsection{Target Ciphers}

We study three lightweight block ciphers, each representing a major design paradigm:

\begin{table}[h]
\centering
\caption{Target cipher specifications.}
\label{tab:ciphers}
\begin{tabular}{@{}lccccl@{}}
\toprule
\textbf{Cipher} & \textbf{Family} & \textbf{Block} & \textbf{Key} & \textbf{Rounds} & \textbf{Input diff.\ $\Delta P$} \\
\midrule
\speck  & ARX     & 32 & 64 & 22 & \texttt{(0x0040, 0x0000)} \\
\simon  & Feistel & 32 & 64 & 32 & \texttt{(0x0000, 0x0001)} \\
\present & SPN    & 64 & 80 & 31 & \texttt{0x0000000000000001} \\
\bottomrule
\end{tabular}
\end{table}

\noindent
\speck uses modular addition, bitwise rotation, and XOR.
\simon uses a balanced Feistel structure with AND, rotation, and XOR.
\present uses a 4-bit S-box layer followed by a bit permutation.
All input differences follow Gohr's convention~\citep{gohr2019} or the natural single-bit choice for the respective cipher.

\subsection{Transfer Protocol}

To measure transfer, we train a distinguisher $f^{(c,r)}$ on cipher $c$ at round $r$ and evaluate it on data from:
\begin{itemize}[itemsep=1pt]
  \item \textbf{Cross-round:} Same cipher, different round count $r' \neq r$.
  \item \textbf{Cross-cipher:} Different cipher $c' \neq c$, same representation.
\end{itemize}
We define \emph{anti-transfer} as accuracy significantly below 50\% ($p < 0.05$) and \emph{positive transfer} as accuracy significantly above 50\%.
The distinction matters: anti-transfer implies the learned features are not merely irrelevant but \emph{actively misleading} on the target distribution.
